\documentclass[11pt]{article}

\usepackage[letterpaper,margin=1in]{geometry}
\usepackage{microtype}
\usepackage{amsmath}
\usepackage{amssymb}
\usepackage{booktabs}
\usepackage{array}
\usepackage[table,dvipsnames]{xcolor}
\usepackage{enumitem}
\usepackage{tikz}
\usetikzlibrary{positioning,arrows.meta,shapes.geometric,fit,backgrounds,calc}
\usepackage[most]{tcolorbox}
\usepackage{float}
\usepackage[hidelinks]{hyperref}

\emergencystretch=3em  % absorb the occasional unbreakable prose line

% Plain-language callout: a non-technical summary readers can rely on without
% reading the math around it.
\newtcolorbox{plain}{
  enhanced, breakable,
  colback=Sepia!7, colframe=Sepia!55, boxrule=0.4pt, arc=2pt,
  left=7pt, right=7pt, top=4pt, bottom=4pt,
  fontupper=\small,
  coltitle=Sepia!72!black, fonttitle=\bfseries\sffamily\footnotesize,
  title={IN PLAIN TERMS}
}

% Definition callout: for normative-ish local definitions used by the method.
\newtcolorbox{definitionbox}{
  enhanced, breakable,
  colback=infoblue!4, colframe=infoblue!45, boxrule=0.4pt, arc=2pt,
  left=7pt, right=7pt, top=4pt, bottom=4pt,
  fontupper=\small,
  coltitle=infoblue!80!black, fonttitle=\bfseries\sffamily\footnotesize,
  title={DEFINITION}
}

% Non-normative aside: implementation guidance and related tooling that inform,
% but are not part of, the normative method.
\newtcolorbox{guidance}{
  enhanced, breakable,
  colback=low!5, colframe=low!45, boxrule=0.4pt, arc=2pt,
  left=7pt, right=7pt, top=4pt, bottom=4pt,
  fontupper=\small,
  coltitle=low!45!black, fonttitle=\bfseries\sffamily\footnotesize,
  title={IMPLEMENTATION GUIDANCE (NON-NORMATIVE)}
}

% --- palette ---
\definecolor{crit}{HTML}{C0392B}
\definecolor{high}{HTML}{E67E22}
\definecolor{med}{HTML}{F1C40F}
\definecolor{low}{HTML}{27AE60}
\definecolor{infoblue}{HTML}{2C3E50}
\definecolor{tagbg}{HTML}{ECF0F1}
\definecolor{rule}{HTML}{BDC3C7}

\newcommand{\tagn}[1]{{\footnotesize\ttfamily #1}}
\newcommand{\code}[1]{\texttt{#1}}
% Category identifiers carry underscores; print them literally in \texttt and
% allow a line break after each underscore so long names never run into the margin.
\makeatletter
\DeclareRobustCommand{\cat}[1]{\texttt{\cat@scan#1_\cat@end}}
\def\cat@scan#1_#2\cat@end{%
  #1%
  \def\cat@rest{#2}%
  \ifx\cat@rest\@empty\else\_\allowbreak\cat@scan#2\cat@end\fi}
\makeatother
\newcommand{\rfckw}[1]{\textsc{#1}}

% Pinned FedRAMP sources, at the same reviewed commit as the companion CVSS memo.
\newcommand{\fedbase}{https://github.com/FedRAMP/2026-markdown/blob/2a75592f0e1f86a48b2888d29af6c570c419a7a5}
\newcommand{\defvulncite}{\href{\fedbase/definitions.md\#vulnerability}{FedRAMP/2026-markdown@2a75592}}
\newcommand{\deflevcite}{\href{\fedbase/definitions.md\#likely-exploitable-vulnerability-lev}{FedRAMP/2026-markdown@2a75592}}
\newcommand{\verpaincite}{\href{\fedbase/providers/rev5/rules/vulnerability-evaluation-and-reporting.md\#estimate-potential-agency-impact}{FedRAMP/2026-markdown@2a75592}}
\newcommand{\verfactcite}{\href{\fedbase/providers/rev5/rules/vulnerability-evaluation-and-reporting.md\#evaluation-factors}{FedRAMP/2026-markdown@2a75592}}
\newcommand{\vdrmitcite}{\href{\fedbase/providers/rev5/rules/vulnerability-detection-and-response.md\#mitigation-and-remediation-expectations}{FedRAMP/2026-markdown@2a75592}}
\newcommand{\ciscite}{\href{https://www.cisecurity.org/cis-benchmarks/cis-benchmarks-faq}{cisecurity.org/cis-benchmarks-faq}}
\newcommand{\scccite}{\href{https://docs.cloud.google.com/security-command-center/docs/attack-exposure-learn}{docs.cloud.google.com/security-command-center/docs/attack-exposure-learn}}

\title{\bfseries A Deterministic PAIN and Remediation Method for\\
FedRAMP Compliance and Security-Benchmark Findings}
\author{Matthew Venne \\ \small Chief Technology Officer, stackArmor}
\date{July 2026}

\begin{document}
\maketitle

\begin{abstract}
\noindent
FedRAMP's Vulnerability Detection and Response (VDR) and Vulnerability Evaluation
and Reporting (VER) rules treat a vulnerability as any weakness --- not only a CVE
but a misconfiguration, an exposure, or a control gap. Security-benchmark scanners
(CIS, DISA STIG, cloud-configuration checks) therefore emit FedRAMP vulnerabilities,
yet they rarely emit the inputs the companion CVSS-Environmental method needs: a
benchmark check reports a rule identifier and a Pass/Fail, seldom a C/I/A impact
vector, and often no severity at all. This memo supplies a deterministic classifier
for that data. A severity-by-asset effect matrix (mirroring the CVE method's
impact-by-requirements combination) produces a customer-effect word; FedRAMP's
effect-and-scope table turns it into a PAIN level; a finding-level
\emph{internet-exercisability} test selects the remediation column; and the same
pinned \code{VDR-TFR-PVR} matrix used by the companion memo selects the deadline.
The method consumes only what scanners and ordinary asset inventory already provide,
fails safe and loud on provider-controlled unknowns, and discloses its one
calibrated default rather than hiding it. A back-test on a real, sanitized CIS Google
Cloud Platform scan of 1{,}378 findings shows no incident-class flood: fourteen
findings (1.0\%) ride the fast clock. This is a companion to \emph{A Deterministic,
CVSS--Environmental Method for FedRAMP Rev5 VDR/VER Vulnerability Prioritization};
the two share notation, the asset model, and the deadline matrix, and are meant to
be read as one family.
\end{abstract}

\tableofcontents

\section{Status of this memo and terminology}
This document is informational. It does not define a FedRAMP requirement; it
specifies a standardized \emph{method} for satisfying the existing VDR/VER
requirements when the finding under evaluation is a security-benchmark or
compliance result rather than a CVE.\footnote{The VDR and VER rules launched on
2026-06-24 as part of the FedRAMP Consolidated Rules for 2026 and apply to both
FedRAMP 20x and Rev5 offerings.} It is a companion to the CVSS-Environmental memo
(hereafter ``the CVE memo''); it reuses that memo's asset model, its remediation
matrix, and its treatment of dispositions, and cites it rather than restating it.

The key words \rfckw{must}, \rfckw{must not}, \rfckw{should}, \rfckw{should not},
and \rfckw{may} are to be interpreted as described in RFC~2119.

\begin{description}[leftmargin=1.5em,style=nextline]
\item[PAIN] Potential Agency Impact, the FedRAMP N1--N5 rating of a finding.
\item[Grade] This method's governed severity of a benchmark finding: Major,
  Moderate, or Minor.
\item[Customer effect] The FedRAMP effect word a finding lands on: Minimal, Narrow,
  Disruptive, or Debilitating.
\item[CR/IR/AR] The CVSS Environmental Confidentiality / Integrity / Availability
  \emph{Requirements} of the affected asset; the asset band summarizes them.
\item[LEV / IRV] Likely Exploitable / Internet-Reachable Vulnerability, the two VER
  axes that select the remediation column.
\item[Internet-exercisable] This method's operational test for whether an
  unauthenticated internet actor can exercise a benchmark finding as it stands. It
  is defined in Section~\ref{sec:exerc} and is \emph{not} a redefinition of IRV.
\item[Class] The provider Certification Class (A--D) selecting a remediation table.
\end{description}

Three terms are held strictly apart throughout and are never interchangeable.
\emph{Internet-accessible} names an entry point --- a listener an outside actor can
reach directly. \emph{Internet-reachable} is FedRAMP's broader status: a finding
that might be triggered by a payload originating on the public internet, including
by an indirect path with no direct route. \emph{Internet-exercisable} is this
method's own operational test for benchmark findings, defined below; it is a
deliberately conservative proxy, not a synonym for either of the other two.

\section{Scope: benchmark findings are FedRAMP vulnerabilities}
\label{sec:scope}
\begin{plain}
A vulnerability, to FedRAMP, is any weakness an attacker could use --- not just a
numbered CVE. A misconfigured storage bucket, a disabled log, an open management
port, a missing encryption setting: each is a vulnerability the rules expect a
provider to detect, evaluate, and remediate on a clock. The trouble is that the
scanners that find these things speak a different dialect from CVE scanners. They
tell you which rule failed and on what resource, but usually not ``how bad'' in the
CVE sense and often not even a severity. This memo builds the missing bridge from
that thin data to a defensible PAIN level and deadline.
\end{plain}

FedRAMP's definition of a vulnerability is deliberately broad: it covers control
gaps, misconfigurations, exposures, and other weaknesses, not only
CVEs.\footnote{FedRAMP, \emph{Definitions} (\emph{vulnerability}). Pinned permalink:
\defvulncite{} (commit of 2026-06-24); accessed 2026-07-16.} A security-benchmark
failure is therefore a FedRAMP vulnerability, and the VDR/VER duties --- estimate
its Potential Agency Impact, evaluate exploitability and reachability, and remediate
within the timeframe matrix --- apply to it in full. The CVE memo answers those
duties for findings that arrive with a CVSS vector. Benchmark findings do not, and
that gap is what this memo closes.

The obstacle is the input data, not the output vocabulary. A CVE scanner emits an
impact vector the CVSS-Environmental formula can consume directly. A benchmark
scanner emits a rule or category identifier and a Pass/Fail state. Some add a
severity of their own, but those labels are not uniform across tools. CIS Benchmarks
publish no severity at all: the Level~1 and Level~2 profiles are hardening
\emph{scopes}, not severities, and \rfckw{must not} be read as one.\footnote{CIS,
\emph{CIS Benchmarks FAQ} (Level 1 and Level 2 profiles). \ciscite{}; accessed
2026-07-16.} DISA STIG content is the exception that proves the rule: it carries a
severity category (CAT~I/II/III) that scanners translate faithfully. So the method
must work from a Pass/Fail and a category name, respect a governed severity when one
genuinely exists, and never invent one where none does.

Two design commitments follow, and both are load-bearing. First, the method reuses
the asset metadata the CVE method already requires --- the affected asset's CR/IR/AR
(or an archetype that resolves to them), its multi-agency scope, and the offering's
Certification Class --- and asks for no benchmark-specific metadata that ordinary
scanners and inventory do not already provide. The only governed artifacts it adds
are a severity source registry and a small per-benchmark determination table, both
described below. Second, exploitability and reachability stay on the axes FedRAMP
built for them: they select the remediation \emph{column}, never the PAIN level. A
single fact is never counted twice.

The normative FedRAMP anchors for what follows are the VER estimate-impact
clause,\footnote{FedRAMP, \emph{VER} (\emph{Estimate Potential Agency Impact}):
\verpaincite{}.} the VER evaluation factors,\footnote{FedRAMP, \emph{VER}
(\emph{Evaluation Factors}): \verfactcite{}.} the VDR mitigation-and-remediation
expectations,\footnote{FedRAMP, \emph{VDR} (\emph{Mitigation and Remediation
Expectations}): \vdrmitcite{}.} and the LEV definition,\footnote{FedRAMP,
\emph{Definitions} (\emph{Likely Exploitable Vulnerability}): \deflevcite{}.} each
pinned at commit 2a75592 (2026-06-24).

\section{The disposition gate}
\label{sec:gate}
\begin{plain}
Before a finding is scored, it has to be real. A scanner's hit can be a false alarm,
a check that does not apply to this system's role, or a weakness already neutralized
by something else. Sorting those out is the disposition step, and it happens
per-finding with evidence --- not by quietly softening the scoring formula to
tolerate noise. The scoring engine assumes what it is handed is a genuine failure;
this gate is what makes that assumption safe.
\end{plain}

Only a confirmed failed condition enters the PAIN calculation. Every other outcome
is a disposition, recorded with its evidence:

\begin{itemize}[leftmargin=1.4em,nosep,topsep=3pt]
\item A \code{Pass} is not a finding.
\item A result that does not match the effective state is a \textbf{false positive}.
\item A recommendation that does not apply to the asset's actual role is
  \textbf{not applicable}. An N/A disposition is an auditable out-of-scope decision
  and carries a mandatory reportable rationale --- it is not a silent drop.
\item A failed configuration neutralized by an effective control is
  \textbf{fully mitigated}; the evidence and the retained PAIN are recorded, exactly
  as in the CVE memo's disposition overlay.
\item A weakness intentionally justified rather than eliminated is an
  \textbf{accepted} finding. Acceptance does not lower the computed PAIN.
\item A weakness reduced but not eliminated by a compensating control is
  \textbf{partially mitigated}: the grade and effect are re-evaluated against the
  \emph{residual} condition on cited evidence, and both the pre- and
  post-mitigation N-ratings are recorded. This matches VER's expectation that a
  remediation clock is satisfied by partial mitigation to a lower N-rating.
\end{itemize}

The scoring algorithm is never weakened to compensate globally for scanner noise.
Noise is resolved here, one finding at a time, against evidence. That discipline is
what lets the rest of the method stay deterministic.

\section{The asset-value band}
\label{sec:band}
\begin{plain}
The same weakness matters more on a crown-jewel system than on a scratch sandbox, so
the first thing the method needs is a sense of how much the affected asset is worth.
It reuses the CVE method's answer exactly: each asset already carries three
criticality dials --- for secrecy, correctness, and uptime. We collapse those three
into one word --- High, Medium, or Low --- and that word is the asset half of the
score. If an asset has never been classified, it is treated as top-value until
someone proves otherwise.
\end{plain}

The affected asset's Security Requirements (CR, IR, AR) are the CVSS Environmental
multipliers used throughout the CVE memo: Low${}=0.5$, Medium${}=1.0$,
High${}=1.5$. The method summarizes them into a single band from their mean:
\[
\text{asset\_mean} = \tfrac{1}{3}(\mathrm{CR}+\mathrm{IR}+\mathrm{AR}),\qquad
\text{band} =
\begin{cases}
\text{Low} & \text{asset\_mean} < 0.75\\
\text{Medium} & 0.75 \le \text{asset\_mean} < 1.25\\
\text{High} & \text{asset\_mean} \ge 1.25
\end{cases}
\]
The attainable means are sixths ($0.5,\,0.667,\,\dots,\,1.5$), so no threshold is
ever hit exactly. The formula is the derivation; the rule an operator actually needs
to remember is plainer:

\begin{definitionbox}
\textbf{The band in plain language.} An asset is \textbf{High} when at least two of
its three requirements are High and none is Low; \textbf{Low} when at least two are
Low and none is High; \textbf{Medium} in every other case. A single Low requirement
demotes two Highs to Medium --- \code{HHL} is Medium, not High.
\end{definitionbox}

Mixed vectors like \code{HHL} are uncommon in practice. The three requirements tend
to correlate --- an asset with genuinely High confidentiality and integrity but Low
availability requirements is unusual --- so the demotion is a mathematical property of
the mean rather than an expected operational case. Where a real asset is genuinely
served poorly by the banding, direct CR/IR/AR assignment and affected-resource
substitution (below) cover it: the band is a convenience for the common, correlated
case, not the only path to a score.

Missing or unresolved asset metadata resolves to \code{HHH}${}\rightarrow{}$High.
Asset classification is provider-controlled, so its absence fails safe and loud: an
unclassified asset scores at the top of the band and surfaces itself for
classification rather than hiding. This is the same fail-safe doctrine the CVE memo
applies; the addendum inherits it unchanged.

\paragraph{Affected-resource substitution.} A benchmark finding is often detected at
one resource but really endangers a higher-value one. A finding that reasonably
affects a higher-value dependent resource, an administrative plane, a shared
credential store, backup or logging infrastructure with reach into higher-value
assets, or a control plane is evaluated \emph{against that higher-value resource},
not against the point of detection. Substitution is governed by that named trigger
list so it is not an open-ended judgment. It carries a firm audit obligation: every
substitution records the substituted resource and the rationale, and the decision is
recorded \emph{both} when it is applied \emph{and} when it is declined
(Section~\ref{sec:audit}). A full asset-adjacency graph would make substitution
table-driven; it is noted as an optional refinement, not a baseline input, because
the baseline may not assume metadata ordinary inventory does not provide.

\begin{plain}
\textbf{Substitution in one example.} A STIG finding lands on a bastion host. The
bastion stores no customer data and is classified L/L/L --- Low on every dimension ---
so on its own nameplate it would band Low. But the bastion is the sole administrative
path to the production database, which is classified H/H/H. The finding is scored
against the \emph{database's} High band, not the bastion's Low band, because what the
weakness actually endangers is the database behind it. The audit record captures the
substituted resource (the database) and why (the bastion is its only admin path). And
if a reviewer looks at the same finding and \emph{declines} to substitute --- judging
the weakness genuinely confined to the bastion --- that decision is recorded too, with
its basis, so an assessor can challenge either call rather than only the ones that
moved the score.
\end{plain}

\section{Finding grade and provenance}
\label{sec:grade}
\begin{plain}
Next the method needs to know how serious the failed check is on its own terms. Where
a trustworthy severity exists --- a DISA STIG category, or a scanner whose severity
scheme is documented and stable --- it is used. Where none exists, as with raw CIS
Pass/Fail, the finding takes a middling ``Moderate'' by default. That default is a
deliberate, disclosed calibration, not a shrug: it reflects that most benchmark
checks really are middling, and that stamping hundreds of hardening gaps as
top-severity would bury the few that are genuinely dangerous.
\end{plain}

The method assigns each finding one of three grades:

\begin{table}[H]
\centering
\renewcommand{\arraystretch}{1.25}
\begin{tabular}{>{\raggedright\arraybackslash}p{0.18\textwidth} >{\raggedright\arraybackslash}p{0.70\textwidth}}
\toprule
\textbf{Grade} & \textbf{Sources} \\
\midrule
\textbf{Major} & STIG CAT~I; governed scanner Critical/High \\
\textbf{Moderate} & STIG CAT~II; governed scanner Medium; \textbf{unscored \code{Fail}} (the calibrated prior) \\
\textbf{Minor} & STIG CAT~III; governed scanner Low \\
\bottomrule
\end{tabular}
\caption{Finding grade by governed severity source.}
\label{tab:grade}
\end{table}

Provenance is resolved through a pinned, versioned \textbf{source registry}, and this
is what keeps the grade from being provider-gameable. The registry records, per
benchmark source, whether it is \emph{structurally unscored} (CIS --- eligible for
the Moderate prior) or carries a \emph{governed mapping} (with the scanner's full
severity vocabulary at a pinned version). A scanner on the published governed list is
governed by default and cannot be un-elected. DISA STIG categories are authoritative.
The starter governed list is DISA STIG CATs, Google Security Command Center category
severities, and AWS Foundational Security Best Practices control severities.

Two rules make the direction of every default explicit:

\begin{itemize}[leftmargin=1.4em,nosep,topsep=3pt]
\item A blank or unrecognized severity \emph{from a scored, governed source} fails
  loud to \textbf{Major}. A provider cannot suppress a governed High and quietly ride
  the Moderate prior --- stripping a governed severity fails \emph{up}, not down.
\item The Moderate prior applies \emph{only} to registry-marked structurally-unscored
  sources. A provider can neither self-escalate nor self-reduce a governed severity,
  and CIS Level~1/2 profiles are never read as severities.
\end{itemize}

\paragraph{Why the prior is Moderate, in full.} The default deserves its complete
justification on the record, because a worst-case default would seem safer and is
not. Missing benchmark severity is \emph{structural}, not a provider omission: CIS
publishes no severities at all, so treating their absence as a fail-safe worst case
would be miscalibration rather than caution --- it punishes a property of the
benchmark, not a gap in the provider's evidence. This is the sharp line between the
two kinds of unknown the method handles: a \emph{provider-controlled} unknown (an
unclassified asset, an unknown firewall state, an unknown agency scope) fails safe
and loud, but a \emph{structurally absent} input (benchmark severity) takes the
calibrated prior. Beyond that principle, two empirical facts support Moderate
specifically. CAT~II is the mode of real benchmark content --- most checks are
genuinely middle-weight hardening items. And labeling hundreds of ordinary hardening
deviations ``potentially debilitating multi-agency events'' would destroy the
incident signal for the findings that actually deserve it. A default that cries wolf
on everything protects nothing.

\paragraph{The escalation override.} The prior can be wrong in the minority, and the
method gives that minority a governed home rather than reopening the flood.
Documented evidence that a specific finding yields direct High C/I/A impact on the
(possibly substituted) resource promotes its grade one step
(Moderate${}\rightarrow{}$Major). The step is audited, fenced by the same registry
and assessor-review governance as the source registry, and expected to be rare. An
encryption-at-rest CAT~II on a shared multi-agency datastore is the canonical case:
genuinely under-claimed at its default grade, and now escalable on evidence.

\section{Customer effect and PAIN}
\label{sec:pain}
\begin{plain}
Now the two halves meet. The grade (how serious the check is) and the asset band (how
much the system is worth) combine in a small lookup table to produce a plain-English
effect word. That word, together with whether the system serves one agency or several,
selects the N-level. The asset is a ceiling: a serious check on a throwaway box cannot
climb higher than the box can support. And one boundary is drawn on purpose --- the
line between an ``incident'' and ``maintenance'' sits exactly where it does to keep the
routine flood out of the incident lane.
\end{plain}

The grade and the asset band select a customer-effect word:

\begin{table}[H]
\centering
\renewcommand{\arraystretch}{1.3}
\begin{tabular}{l ccc}
\toprule
\textbf{Grade} & \textbf{Asset High} & \textbf{Asset Medium} & \textbf{Asset Low} \\
\midrule
\textbf{Major}    & \cellcolor{crit!20}Debilitating & \cellcolor{high!25}Disruptive & \cellcolor{med!30}Narrow \\
\textbf{Moderate} & \cellcolor{med!30}Narrow         & \cellcolor{med!30}Narrow      & \cellcolor{low!25}Minimal \\
\textbf{Minor}    & \cellcolor{low!25}Minimal        & \cellcolor{low!25}Minimal     & \cellcolor{low!25}Minimal \\
\bottomrule
\end{tabular}
\caption{Customer-effect matrix (grade $\times$ asset band).}
\label{tab:effect}
\end{table}

The effect word and the asset's agency scope select the PAIN level, using FedRAMP's
fixed effect-and-scope semantics:

\begin{table}[H]
\centering
\renewcommand{\arraystretch}{1.3}
\begin{tabular}{l cc}
\toprule
\textbf{Customer effect} & \textbf{Single agency} & \textbf{Multi-agency} \\
\midrule
Debilitating & N4 & \cellcolor{crit!20}N5 \\
Disruptive   & N3 & \cellcolor{high!25}N4 \\
Narrow       & N2 & N2 \\
Minimal      & N1 & N1 \\
\bottomrule
\end{tabular}
\caption{Effect $\times$ agency scope $\rightarrow$ PAIN. N5 is reachable only for a
Debilitating multi-agency finding.}
\label{tab:paintbl}
\end{table}

Unknown agency scope resolves to multi-agency: scope is provider-controlled metadata,
so its absence fails safe and loud. The band acts as a ceiling in
Table~\ref{tab:effect}: severity never raises effect above what the asset can
support, so a Major finding on a Low-band asset is still only Narrow. And the
Moderate${}\rightarrow{}$Minor grade transition never drops the effect by more than
one band (exactly one, unless the effect is already Minimal), which keeps the matrix
monotone.

\paragraph{Two observations the tables make plain.} First, \textbf{the multi-agency
column has no N3.} This is FedRAMP's own definition, not an artifact of this method:
N3 is defined as an effect disruptive to a single agency, and disruption of more than
one agency is N4 by definition. No classifier could put a genuinely multi-agency
Disruptive effect on N3 and still be faithful to the rule. The line between an
incident (N4 and above) and maintenance (N2 and below) therefore sits exactly on the
Disruptive boundary, and the effect matrix is calibrated with that boundary in view.

That said, N3 is not unreachable on shared infrastructure --- because the scope input
is the agency scope of the finding's \emph{effect}, not the asset's nameplate tenancy.
A finding on a shared component whose exploitation is demonstrably contained to a
single agency's tenancy --- a per-tenant bucket-prefix policy error, a tenant-scoped
IAM binding that leaks only within one tenant --- scores single-agency, so a Disruptive
such finding lands on N3, not N4. That single-agency assertion on a shared asset is
exactly the case the audit record already requires a tenancy basis for
(Section~\ref{sec:audit}); the containment must be evidenced, not asserted. The
fail-safe is unchanged: unknown scope, or a finding on a shared control plane whose
blast radius is not demonstrably contained, resolves to multi-agency.

Second, \textbf{the Moderate/High cell is a disclosed calibration knob.} It reads
Narrow, not Disruptive. The alternative --- Disruptive --- would put every CAT~II and
every raw CIS \code{Fail} on a high-value multi-agency asset at N4, which is
incident-class in Class~D. That is precisely the flood this design exists to prevent,
so the knob is set to Narrow deliberately and disclosed here rather than buried. The
choice has a named residual risk, and the method owns it rather than papering over it:

\begin{definitionbox}
\textbf{Named residual risk (the N5$\leftrightarrow$N2 cliff on High assets).} Because
Moderate/High caps at Narrow${}\rightarrow{}$N2, a genuinely dangerous CAT~II on a
crown-jewel multi-agency asset --- the standing example is an account-lockout policy
gap on an internet-exercisable admin plane, brute-forceable in practice --- is capped
at N2 unless the escalation override (Section~\ref{sec:grade}) is invoked on evidence.
The gap between that N2 and the N5 the finding might deserve is the calibration's
known cliff. It is not closed by moving the matrix cell (that would let exercisability
set PAIN, the double-count the method forbids); it is closed, when it applies, by the
evidence-gated grade escalation, and it is owned by the certifying provider.
\end{definitionbox}

\section{Internet exercisability: the remediation column}
\label{sec:exerc}
\begin{plain}
The N-level says how bad; this step says how fast. FedRAMP tightens the clock for
findings that are likely exploitable and internet-reachable. For benchmark findings
we cannot measure those directly, so the method uses one honest, conservative proxy:
can an unauthenticated stranger on the internet actually exercise this failed
setting, as it stands, without first getting something the finding itself does not
hand them? If yes, it rides the fast clock. This test is a practical stand-in, not a
new definition of FedRAMP's reachability --- and it only ever errs toward calling
something \emph{not} exercisable, never the other way.
\end{plain}

The normative test is a single sentence: \textbf{a benchmark finding is
internet-exercisable if and only if an unauthenticated internet actor can exercise
the failed condition as it currently stands, with no precondition the finding does
not itself supply.} It is the same spirit as the \code{FRD-LEV}
unauthenticated-automation floor the CVE memo uses for direct-IRV CVEs.

\begin{definitionbox}
Internet-exercisability is an \textbf{operational test for benchmark findings}. It is
\emph{not} a redefinition of FedRAMP's Internet-Reachable Vulnerability status, and it
must not be read as one. The relationship is one-directional and stated as a soundness
claim: \textbf{exercisable ${}\Longrightarrow{}$ IRV, never the converse.} A finding
the test calls exercisable is genuinely internet-reachable; a finding it calls not
exercisable may still be IRV by an indirect path the proxy cannot see. The baseline is
therefore \emph{sound but incomplete} with respect to IRV --- it can only ever
\emph{under}-select the fast clock, never over-select it --- and the evidence-backed
overrides below exist to catch the cases it misses.
\end{definitionbox}

\paragraph{Per-family operational rules.} The test is evaluated by benchmark family,
each family using only data ordinary scanners and inventory provide:

\begin{itemize}[leftmargin=1.4em,topsep=3pt]
\item \textbf{Host benchmarks, OS-level} (RHEL/Windows STIG and CIS): exercisable iff
  the administrative/login plane is world-open, tested against a governed port set
  (default $\{22, 3389, 5985, 5986, 5900, 10250\}$). The \code{admin\_plane\_open}
  bit is not scanner-emitted metadata --- it is the output of a \emph{required join}
  against governed firewall/load-balancer inventory, the same feed the CSP
  benchmark's own network controls read. There is no per-rule mapping. Unknown
  firewall state resolves to open: firewall state is a discoverable unknown, and the
  doctrine is to score discoverable unknowns loud to force the inventory join. A
  WAF/ALB-fronted surface counts as open --- a WAF is mitigation, not closure.
\item \textbf{Host benchmarks, application-level} (web/db STIGs): exercisable iff that
  application's own service surface is internet-open, by the same join.
\item \textbf{Cloud-configuration benchmarks:} exercisable iff the finding category is
  \emph{exposure-class} under the governed determination table below.
\end{itemize}

\paragraph{The determination table.} For cloud benchmarks the normative object is a
governed, content-hashed table published per (benchmark, version, scanner). The table
assigns every category identifier in the version's vocabulary an effective
classification: \emph{exposure-class} (always exercisable), \emph{surface} (exercisable
iff the attached resource is public, joined from the same scan's exposure findings;
unknown attachment resolves to public), or
\emph{identity/operations-plane} (not exercisable). The classification reads the
category identifier the scanner emits --- never the rule prose, and never per-resource
IAM analysis --- so the determination is a lookup, and a second implementer reproduces
any result from exactly two inputs in the audit record: the table hash and the category.
A scanned category \emph{absent} from the table's vocabulary --- a scanner's monthly
new addition, say --- fails loud to LEV+IRV and is flagged unreviewed.

That is the whole normative requirement: the table, its three classes, the fail-loud
rule, and the governance below. \emph{How a canonical table is built} --- what makes a
category exposure-class in the first place --- is implementation guidance:

\begin{guidance}
The recommended generator is token and literal-pattern matching on the identifier.
Normalize --- uppercase, split on non-alphanumeric runs --- and classify exposure-class
if any token is one of $\{$\code{PUBLIC}, \code{OPEN}, \code{ANONYMOUS},
\code{INTERNET}, \code{WORLD}$\}$. Token boundaries make the match safe across scanner
grammars: \cat{OPEN_SSH_PORT} matches on the \code{OPEN} token while
\cat{OPENSSH_CONFIG} does not, and an AWS Config identifier such as
\code{s3-bucket-public-read-prohibited} normalizes to the same tokens as its GCP
cousin. Before tokenizing, match two literal patterns: an identifier containing
\code{0.0.0.0} or \code{::/0} is exposure-class --- caught by literal substring
precisely because normalization would split \code{0.0.0.0} into a bare \code{0} and
lose the meaning. So \code{sg-allow-ingress-from-0.0.0.0-0} matches, while the plain
numeric segments of \cat{TLS_1_0_ENABLED} do not. Because the generator reads
scanner-emitted identifiers, a provider cannot game it by renaming; because its output
is frozen into the reviewed, hashed table, the generator itself never runs at scoring
time. In practice a canonical table is serialized as review deltas against the
generator --- \code{add}, \code{remove}, and \code{surface} entries plus the
\code{vocabulary} --- so review effort concentrates on deviations. A category left
untagged in that serialization is not a gap: it is one the review affirmatively
classified identity/operations-plane.

Two tokens that look tempting --- \code{EXTERNAL} and \code{UNRESTRICTED} --- are
deliberately \emph{not} in the auto-match set. Their genuine true positives are
already carried by \code{PUBLIC}/\code{OPEN}/\code{INTERNET} or by explicit
\code{add} entries, while their false-positive rate is high
(\cat{SQL_EXTERNAL_SCRIPTS_ENABLED} is not an internet exposure;
\cat{API_KEY_APIS_UNRESTRICTED} scopes a key, not a surface). Keeping them out holds
the auto-match vocabulary tight and moves the burden of proof onto a reviewed table
entry. Screening a benchmark's full category list for token collisions when the table
is built is a sound practice, not a normative requirement. So is richer analysis:
automated or agentic analysis of a benchmark's category list to pre-classify exposure
relevance for triage is permitted and useful, but it is not required --- token and
literal-pattern matching against a reviewed table meets the bar.
\end{guidance}

\paragraph{Governance.} The governance around the table is what makes it trustworthy
rather than a hidden lever:

\begin{itemize}[leftmargin=1.4em,nosep,topsep=3pt]
\item Canonical tables are published per (benchmark, version, scanner) and
  \textbf{content-hashed}; the hash travels in the audit record. A version
  \emph{string} enforces nothing; a hash does.
\item Any provider deviation from the canonical table is enumerated per-entry with a
  rationale and flagged for assessor review.
\item Reclassifying a category out of exposure-class is a hard assessor-review
  trigger --- it is the one move that silently slows a clock, so it never happens
  silently.
\item Keeping the tables current is an operating requirement with a monthly
  new-category diff, not a once-per-version event.
\end{itemize}

The identity/operations-plane class deserves its justification stated once, because
it is where most of a benchmark lands. The test's own precondition clause supplies
it: a user-managed service-account key requires possessing the secret, and the public
cloud control plane alone gives an internet actor nothing to exercise. A data-plane
grant to \code{allUsers} or \code{allAuthenticatedUsers}, by contrast, \emph{is}
exposure-class regardless of the category name --- \code{allAuthenticatedUsers} is
satisfied by any self-registered account, which the unauthenticated test treats as
self-suppliable.

\paragraph{Overrides.} The overrides carry the cases the sound-but-incomplete
baseline misses, and each is evidence-backed, audited, and expected to be uncommon:

\begin{itemize}[leftmargin=1.4em,nosep,topsep=3pt]
\item \textbf{LEV+NIRV} when evidence shows a likely tenant, internal, adjacent, or
  local actor can exercise the condition. The canonical case is the MFA-not-enforced
  family (Section~\ref{sec:examples}).
\item \textbf{LEV+IRV} when evidence shows public usability (for example a
  known-leaked key, at which point incident response applies), or on
  transitive-reachability evidence --- an indirect internet payload path, mirroring
  the CVE companion's transitive branch --- which gives the VER indirect-payload case
  a home in the IRV direction.
\end{itemize}

The column then selects itself: LEV+IRV if internet-exercisable or so overridden;
LEV+NIRV on an evidence-backed internal-exploitability override; NLEV otherwise. The
determination is stated once more for emphasis: \textbf{exercisability changes the
column only; it never changes PAIN.} An affirmative not-exercisable determination
(any \code{False} on \code{admin\_plane\_open}, \code{service\_open}, or
\code{attached\_resource\_public}) requires an enforced-state evidence pointer --- a
security-group or firewall rule identifier, a listener binding --- because the absence
of an observation is not the same as enforced closure.

\section{Deadlines}
\label{sec:deadline}
\begin{plain}
The last step is a table lookup. The PAIN level, the remediation column, and the
provider's assurance Class pick a single cell, and that cell is the number of days
allowed. The numbers are FedRAMP's, not ours --- the very same matrix the CVE method
uses --- so nothing benchmark-specific is invented here.
\end{plain}

The deadline is
$\text{deadline} = M[\text{Class}][\text{PAIN}][\text{column}]$, drawn from the pinned
\code{VDR-TFR-PVR} matrix reproduced in the CVE memo's appendix on remediation
deadline matrices (days; $0.5 = 12$ hours; Classes~A and~B share a schedule; N1
carries no FedRAMP deadline). The full tables are not restated here; only the single
N2 row the worked examples use is reproduced for convenience:

\begin{table}[H]
\centering
\footnotesize
\renewcommand{\arraystretch}{1.2}
\begin{tabular}{l ccc @{\hspace{2em}} ccc @{\hspace{2em}} ccc}
\toprule
& \multicolumn{3}{c}{\textbf{Class A / B}} & \multicolumn{3}{c}{\textbf{Class C}} & \multicolumn{3}{c}{\textbf{Class D}} \\
\cmidrule(lr){2-4}\cmidrule(lr){5-7}\cmidrule(lr){8-10}
\textbf{PAIN} & L+I & L+N & NLEV & L+I & L+N & NLEV & L+I & L+N & NLEV \\
\midrule
N2 & 96 & 160 & 192 & 48 & 128 & 192 & 24 & 96 & 192 \\
\bottomrule
\end{tabular}
\caption{The single \code{VDR-TFR-PVR} row used by the worked examples, reproduced from
the CVE memo's appendix. The complete N5--N2 matrix lives there and is not duplicated.}
\label{tab:n2row}
\end{table}

The anchor case ties the whole chain together. A CAT~II on a medium-value,
internet-exercisable host is Moderate${}\times{}$Medium${}\rightarrow{}$Narrow
${}\rightarrow{}$N2, LEV+IRV, which is \textbf{24 days in Class~D} (48 in Class~C, 96
in Classes~A/B) --- and 192 days if it were not exercisable. Exercisability moved the
clock by an order of magnitude and left the N-level untouched.

\paragraph{Where these clocks are actually won.} A remediation clock is the wrong
place to fight compliance-finding volume, and the deadline machinery should not be
mistaken for the strategy. Compliance findings overwhelmingly open at build and deploy
time, not part-way through operations and maintenance: a bucket is created public, a
firewall is opened, a role is over-granted. The practical lever for N4/N5 volume is
therefore \emph{pre-deployment} --- infrastructure-as-code and policy analysis that
evaluates the same two inputs this method needs (asset archetype $\times$ finding
category) before the resource exists. A public bucket is an acceptable pattern in some
designs; a public bucket bound to a data-sensitive archetype is precisely the thing to
block at deploy time rather than remediate on a 12-hour clock afterward. The timeframe
matrix runs from completed evaluation and governs what escapes into production;
prevention is what keeps the incident-class population near zero in the first place.
The back-test that follows should be read in that light --- a small fast-clock count
is the goal, and shifting the check left is how it stays small.

\section{The minimum audit record}
\label{sec:audit}
\begin{plain}
Everything the method decides has to be reconstructable later from what it wrote down.
Two reviewers handed the same finding, the same pinned tables, and the same evidence
must reach the same answer --- so the audit record captures not just the result but
every input and every judgment call, including the ones where the method chose
\emph{not} to act.
\end{plain}

Every output is reproducible from a single record. It captures: the benchmark and
version; the rule or category identifier and the scanner result; the disposition and
its evidence (for anything other than an unqualified \code{Fail}); the affected asset
and its resolved archetype; the affected-resource substitution and rationale when
applied, \emph{and} the substitution decision and basis when declined; the CR/IR/AR
and derived band; the agency scope, with a ``defaulted'' marker if it was unknown; a
tenancy basis (agencies served, plus evidence) for any single-agency assertion on a
shared-service archetype; the severity and its provenance (the grade and why the
source is governed or unscored); the customer effect and PAIN; the exercisability
determination together with the determination-table content hash and the enumerated
provider delta; an enforced-state evidence pointer for any affirmative
not-exercisable determination; the column and any override evidence reference; and the
Certification Class, the evaluation completion time, and the resulting deadline.

\section{Worked examples}
\label{sec:examples}
\begin{plain}
The method run by hand on the cases that decide whether it is calibrated right ---
including the two it deliberately handles differently from what a first glance
suggests (the service-account key, and MFA).
\end{plain}

\paragraph{Example 1 --- user-managed service-account key
(\cat{USER_MANAGED_SERVICE_ACCOUNT_KEY}).} Present in the vocabulary with no exposure
or surface classification: identity/operations-plane
${}\rightarrow{}$ \textbf{NLEV}. The public GCP API plane is not the finding's
exposure; an internet actor lacks the secret, so the unauthenticated test fails at its
precondition. Overrides remain available: NIRV on evidence of broad internal
availability plus harmful permissions, IRV only on evidence of public disclosure.

\emph{Contrast --- API keys (\cat{API_KEY_APPS_UNRESTRICTED}).} The same
possess-the-secret reasoning does \emph{not} transfer. Client-embedded GCP API keys
(Maps, Firebase web and mobile) are public by design; the finding is precisely that
the compensating restriction on an already-public key is missing. These categories are
therefore \code{surface}, not \code{remove}: the attachment is ``key is
client-distributed,'' unknown distribution defaults to exercisable, and a down-override
to NLEV requires evidence the key is server-side-only and network-scoped. Two findings
that both contain \code{UNRESTRICTED} split on their actual exposure, not on a shared
word.

\paragraph{Example 2 --- public bucket (\cat{PUBLIC_BUCKET_ACL}).} Exposure-class in
the determination table (matched on the \code{PUBLIC} token by the generator)
${}\rightarrow{}$ \textbf{LEV+IRV}. Governed SCC severity
High ${}\rightarrow{}$ Major; on a High-value multi-agency data archetype the effect
is Debilitating ${}\rightarrow{}$ \textbf{N5} ${}\rightarrow{}$ \textbf{12 hours in
Class~D}. That is an incident, and the method calls it one.

\paragraph{Example 3 --- \cat{SSL_NOT_ENFORCED} on a Cloud SQL instance where
\cat{SQL_PUBLIC_IP} is Compliant.} A \code{surface} category: exercisable only if the
attached resource is public. The same scan shows the instance has no public IP, so the
attachment joins to False ${}\rightarrow{}$ \textbf{NLEV}. This is the permitted
category-level inference --- when the attaching exposure control is fully Compliant in
the same scan, its resources are non-public by evidence, not assumption.

\paragraph{Example 4 --- CAT~II on a public web VM with admin ports closed.} An
OS-level host finding. The VM serves the public over 443, but the administrative plane
is closed by the firewall join, so the finding is not exercisable via 443
${}\rightarrow{}$ \textbf{NLEV}. On a Medium-band asset it is Moderate${}\times{}$
Medium ${}\rightarrow{}$ Narrow ${}\rightarrow{}$ N2 ${}\rightarrow{}$ \textbf{192
days in Class~D}. The application being internet-facing does not by itself make an
OS-hardening deviation internet-exercisable.

\paragraph{Example 5 --- MFA not enforced.} The strict test says not exercisable: an
internet actor still needs a valid credential, which the finding does not supply. This
is the canonical override case named in Section~\ref{sec:exerc}. A provider facing a
likely tenant or internal actor \rfckw{should} consider the LEV+NIRV override (or a
deliberate LEV+IRV where credential stuffing against a public sign-in is the realistic
threat), recording the evidence once as a standing family rule rather than
finding-by-finding.

\section{Back-test results}
\label{sec:backtest}
\begin{plain}
The real test of a calibration is whether it floods the incident queue on a genuine
scan. It does not. Run against a real, sanitized CIS Google Cloud scan of nearly
fourteen hundred findings --- and under the harshest asset assumption, treating every
asset as top-value --- only fourteen findings land on a fast clock, and the incident
lane stays essentially empty. The bulk sits at the low, slow end, which is exactly
where routine hardening gaps belong.
\end{plain}

The method was back-tested against a sanitized real CIS Google Cloud Platform
Foundation~2.0 scan of \textbf{1{,}378 findings} (Class~D, multi-agency, from Google
Security Command Center). To stress the calibration, the run assumes the
\emph{High} asset band uniformly --- the loudest assumption available. Even so, the
distribution is dominated by the slow, low end:

\begin{table}[H]
\centering
\renewcommand{\arraystretch}{1.25}
\footnotesize
\begin{tabular}{l l r r}
\toprule
\textbf{PAIN} & \textbf{Column} & \textbf{Findings} & \textbf{Class-D deadline} \\
\midrule
N5 & LEV+IRV & 5   & 12 hours \\
N5 & NLEV    & 2   & 8 days \\
N2 & LEV+IRV & 9   & 24 days \\
N2 & NLEV    & 736 & 192 days \\
N1 & NLEV    & 626 & no deadline \\
\midrule
\multicolumn{2}{l}{\textbf{Total}} & \textbf{1{,}378} & \\
\bottomrule
\end{tabular}
\caption{Back-test distribution on the sanitized CIS GCP 2.0 scan (Class~D,
multi-agency, uniform High-band assumption). Only exposure-class findings ride the
LEV+IRV fast clock.}
\label{tab:backtest}
\end{table}

Two properties matter. First, there is \textbf{no N4/N5 flood}: only seven findings
reach incident-class at all, and just five ride the 12-hour clock, on the loudest
possible asset assumption. Second, the fast clock is scarce and earned:
$5 + 9 = 14$ findings ride LEV+IRV, exactly \textbf{1.0\%} of the population, and every
one of them is an exposure-class category. The remainder --- the 736 N2 and 626 N1
findings --- are the ordinary hardening backlog, correctly parked on the slow clocks
where they belong. Dropping the asset assumption to Medium shifts the incident-class
findings from N5 to N4 (a slightly slower incident clock) and leaves the N2/N1 bulk
untouched, which is the intended behavior: the asset band moves the row, never the
column.

\paragraph{A related signal: SCC attack exposure scores.} It is worth noting where a
commercial tool overlaps this method, because the resemblance is close enough to
mislead. Google Security Command Center's attack exposure
scores\footnote{Google, \emph{Security Command Center: Learn about attack exposure
scores}. \scccite{}; accessed 2026-07-17.} are built from structurally the same two
ingredients this method uses: a resource-value configuration (HIGH${}=10$,
MEDIUM${}=5$, LOW${}=1$, the analogue of the asset band) and an internet-origin
attack-path simulation (the analogue of exercisability). The parallel is real, and it
makes SCC a useful neighbor. It cannot, however, replace this method for FedRAMP, for
four concrete reasons.

\begin{guidance}
SCC attack exposure \textbf{complements but cannot replace} the method here.
\emph{Not reproducible from an audit record:} the scoring algorithm is only partially
published and includes probabilistic attack-path simulation recalculated roughly every
six hours, so the same finding can carry different scores on different days and a
reviewer cannot re-derive a score from what is written down --- the determinism
requirement fails. \emph{No FedRAMP semantics:} it emits neither the effect/scope
words nor the N-levels nor the Class-keyed \code{VDR-TFR-PVR} deadlines the rule
requires. \emph{No NIRV analogue:} its simulations model only internet-origin
attackers, so there is no not-internet-reachable branch --- the internal-actor cases
this method routes to LEV+NIRV have no home in it. \emph{Single-provider, gated:} it
requires org-level Premium activation on one cloud. Its right role is as a strong
\emph{evidence source} for the override lanes --- an attack-path result is exactly the
kind of evidence that supports a LEV+NIRV or LEV+IRV override, or a down-override off
the fast clock --- and as an independent \emph{calibration cross-check} against the
back-test distribution.
\end{guidance}

\section{Known limitations}
\label{sec:limits}
\begin{plain}
The method is an honest proxy, and honest proxies have known blind spots. It can
shout a little too loudly about harmless findings that merely sit on an exposed host,
and it can under-rate a finding an insider could exploit unless someone supplies the
evidence to override it. It is a calibrated triage tool for benchmark findings, not a
substitute for a full exploitability analysis, and it is never presented as one.
\end{plain}

The method's approximations are disclosed, not hidden:

\begin{itemize}[leftmargin=1.4em,topsep=3pt]
\item \textbf{It over-prioritizes some non-exercisable findings on exposed surfaces.}
  The host proxy rides the admin-plane bit for the whole host, so a bastion's entire
  CAT~II set can ride LEV+IRV even though most of those findings are not individually
  internet-exercisable. This over-generation is the named, disclosed error shape of
  the host proxy; a login-plane rule-group refinement is available as an optional
  governed extension, and the reporting layer groups findings so the noise is
  legible.
\item \textbf{It under-prioritizes internally exploitable findings absent an
  override.} Because the baseline exercisability test is sound but incomplete
  (exercisable ${}\Longrightarrow{}$ IRV, never the converse), a finding an insider,
  tenant, or adjacent actor could exploit stays on the slow clock until the
  evidence-backed override is applied. The MFA family and tenant-isolation families
  are the standing examples, handled as family rules rather than per-finding.
\item \textbf{The unclassified-asset and unknown-firewall defaults are loud on
  purpose, and that is itself a flood vector.} Operating without an inventory
  classification feed or a firewall feed will over-score, by design --- the loud
  default exists to force those joins. Inventory classification is an operating
  precondition, not a silent assumption.
\item \textbf{It is not a full LEV evaluation.} The method is a calibrated,
  deterministic triage for benchmark findings from the data scanners actually emit. It
  is never described as equivalent to a full likely-exploitability analysis, and where
  a finding warrants that depth, the override and disposition paths route it there.
\end{itemize}

\begin{plain}
\textbf{What to take away.} The method scores benchmark findings using only what
scanners and ordinary inventory already emit --- a rule identifier, a Pass/Fail, an
asset tag, a firewall feed. Every unknown the \emph{provider} controls --- an
unclassified asset, an unknown firewall state, an unknown agency scope --- scores
loud, so silence never makes a problem look smaller than it is. The one place the
method quietly picks a default is the severity of a benchmark that structurally
publishes none, where an unscored \code{Fail} becomes Moderate; that single calibrated
choice is disclosed here in full, not hidden, and it is the reason a real scan of
1{,}378 findings produces fourteen fast-clock items instead of a wall of false
alarms.
\end{plain}

\appendix
\section{Worked determination table: CIS GCP Foundation 2.0}
\label{app:table}
The determination table below is the canonical artifact for the CIS Google Cloud
Platform Foundation~2.0 benchmark as emitted by Google Security Command Center. It is
what the back-test of Section~\ref{sec:backtest} ran against. In production the table
is content-hashed and the hash is carried in each finding's audit record; the
vocabulary is the benchmark version's complete category list, abbreviated here to the
entries the classification logic turns on.

The recommended generator classifies the pure-exposure categories automatically ---
their identifiers carry an exposure token --- and review confirmed each, so they carry
no delta entry and are always exercisable:

\begin{table}[H]
\centering
\renewcommand{\arraystretch}{1.2}
\footnotesize
\begin{tabular}{l l l}
\toprule
\textbf{Category} & \textbf{Matching token} & \textbf{Result} \\
\midrule
\cat{PUBLIC_BUCKET_ACL}   & \code{PUBLIC} & LEV+IRV \\
\cat{PUBLIC_DATASET}      & \code{PUBLIC} & LEV+IRV \\
\cat{PUBLIC_SQL_INSTANCE} & \code{PUBLIC} & LEV+IRV \\
\cat{KMS_PUBLIC_KEY}      & \code{PUBLIC} & LEV+IRV \\
\cat{OPEN_SSH_PORT}       & \code{OPEN}   & LEV+IRV \\
\cat{OPEN_RDP_PORT}       & \code{OPEN}   & LEV+IRV \\
\bottomrule
\end{tabular}
\caption{Exposure-class categories (generator-matched, confirmed on review).}
\label{tab:kw}
\end{table}

The \code{surface} entries are exercisable only when the attached resource is public,
joined from the same scan's exposure state; unknown attachment resolves to public
(fail-safe toward the fast clock):

\begin{table}[H]
\centering
\renewcommand{\arraystretch}{1.2}
\footnotesize
\begin{tabular}{l >{\raggedright\arraybackslash}p{0.62\textwidth}}
\toprule
\textbf{\code{surface} category} & \textbf{Attachment tested} \\
\midrule
\cat{SSL_NOT_ENFORCED}          & Cloud SQL instance public (via \cat{SQL_PUBLIC_IP}) \\
\cat{WEAK_SSL_POLICY}           & fronting load balancer public \\
\cat{SQL_NO_ROOT_PASSWORD}      & Cloud SQL instance public \\
\cat{PUBLIC_IP_ADDRESS}         & firewall/authorized-network state (reachability, not an exercisable surface) \\
\cat{SQL_PUBLIC_IP}             & firewall/authorized-network state \\
\cat{API_KEY_APPS_UNRESTRICTED} & key is client-distributed (public by design; unknown $\rightarrow$ exercisable) \\
\cat{API_KEY_APIS_UNRESTRICTED} & key is client-distributed \\
\bottomrule
\end{tabular}
\caption{\code{surface} categories: exercisable iff the attached resource is public.}
\label{tab:surface}
\end{table}

The \code{add} and \code{remove} sets are both empty for this version: the generator
covers every pure-exposure category, and after the API-key categories were moved to
\code{surface} there is no generator false positive left to strip. Every other
category in the vocabulary --- the logging, monitoring, IAM, CMEK, and SQL-hardening
families such as \cat{USER_MANAGED_SERVICE_ACCOUNT_KEY}, \cat{MFA_NOT_ENFORCED},
\cat{DNSSEC_DISABLED}, \cat{IP_FORWARDING_ENABLED}, and
\cat{COMPUTE_SERIAL_PORTS_ENABLED} --- is identity/operations-plane and resolves to
NLEV, with the MFA and tenant-isolation families carrying the standing overrides of
Section~\ref{sec:exerc}. Integrity-hardening categories whose exploitation needs an
off-path or race precondition (\cat{DNSSEC_DISABLED}, \cat{RSASHA1_FOR_SIGNING}) fail
the unauthenticated-exercise test and take the NLEV default with an override available.
A category scanned but \emph{absent} from the vocabulary fails loud to LEV+IRV and is
flagged unreviewed.

\section*{Normative sources}
\addcontentsline{toc}{section}{Normative sources}
\begingroup\sloppy
\begin{itemize}[leftmargin=3.4em]
\item[{[DEF]}] FedRAMP, \emph{FedRAMP Definitions}. Definitions cited:
  \emph{vulnerability} (\defvulncite{}), \emph{Likely Exploitable Vulnerability}
  (\deflevcite{}). Commit of 2026-06-24; accessed 2026-07-16.
\item[{[VER]}] FedRAMP, \emph{Vulnerability Evaluation and Reporting (VER)}, Rev5.
  Sections cited: \emph{Estimate Potential Agency Impact} (\verpaincite{}),
  \emph{Evaluation Factors} (\verfactcite{}). Commit of 2026-06-24.
\item[{[VDR]}] FedRAMP, \emph{Vulnerability Detection and Response (VDR)}, Rev5.
  Section cited: \emph{Mitigation and Remediation Expectations} (\vdrmitcite{}); and
  the \code{VDR-TFR-PVR} remediation matrix, reproduced in the CVE memo's appendix.
  Commit of 2026-06-24.
\item[{[CIS]}] CIS, \emph{CIS Benchmarks FAQ} (Level~1 and Level~2 profiles).
  \ciscite{}; accessed 2026-07-16.
\item[{[CVEM]}] M.~Venne, \emph{A Deterministic, CVSS--Environmental Method for
  FedRAMP Rev5 VDR/VER Vulnerability Prioritization}, stackArmor, July 2026. Companion
  memo supplying the asset model, the disposition overlay, and the \code{VDR-TFR-PVR}
  deadline matrix.
\item[{[BOD]}] CISA, \emph{Binding Operational Directive 26-04}, 2026-06-10. The
  mandating authority for the VDR/VER timeframes.
\end{itemize}
\endgroup

\end{document}
